% !TeX program = lualatex
\documentclass[10pt,a4paper]{article}

% ---------------------------------------------------------------------------- %
% PACKAGES
% ---------------------------------------------------------------------------- %

\usepackage[
    top=1.25cm,
    bottom=1.25cm,
    left=1.7cm,
    right=1.7cm
]{geometry}

\usepackage[T1]{fontenc}
\usepackage[utf8]{inputenc}
\usepackage[scaled=0.94]{helvet}
\usepackage{xcolor}
\usepackage{hyperref}
\usepackage{enumitem}
\usepackage{tabularx}
\usepackage{array}

\pagestyle{empty}

% ---------------------------------------------------------------------------- %
% ---------------------------------------------------------------------------- %
% CV CONFIGURATION
% ---------------------------------------------------------------------------- %
% ---------------------------------------------------------------------------- %
%
% All visual formatting should be controlled from this section.
%
% ---------------------------------------------------------------------------- %


% ---------------------------------------------------------------------------- %
% FONT
% ---------------------------------------------------------------------------- %

% Helvetica-compatible sans serif supplied by the standard TeX installation.
% This avoids depending on system fonts or the fontspec font loader.
\renewcommand{\familydefault}{\sfdefault}


% ---------------------------------------------------------------------------- %
% COLOURS
% ---------------------------------------------------------------------------- %

\definecolor{CVBlack}{HTML}{000000}
\definecolor{CVDarkGray}{HTML}{333333}
\definecolor{CVGray}{HTML}{666666}
\definecolor{CVLightGray}{HTML}{CCCCCC}


% ---------------------------------------------------------------------------- %
% FONT SIZES
% ---------------------------------------------------------------------------- %

% ---------- Header ----------

\newcommand{\KeywordSize}{9pt}
\newcommand{\NameSize}{20pt}
\newcommand{\ContactSize}{8pt}

% ---------- Opening statement ----------

\newcommand{\StatementSize}{9.5pt}

% ---------- Section headings ----------

\newcommand{\SectionSize}{11.5pt}

% ---------- CV entries ----------

\newcommand{\RoleSize}{10.3pt}
\newcommand{\CompanySize}{9.6pt}
\newcommand{\DurationSize}{9pt}
\newcommand{\BodySize}{9.2pt}


% ---------------------------------------------------------------------------- %
% LINE HEIGHTS
% ---------------------------------------------------------------------------- %

% ---------- Header ----------

\newcommand{\KeywordLineSize}{0pt}
\newcommand{\NameLineSize}{20pt}
\newcommand{\ContactLineSize}{8pt}

% ---------- Opening statement ----------

\newcommand{\StatementLineSize}{10pt}

% ---------- Section headings ----------

\newcommand{\SectionLineSize}{12pt}

% ---------- CV entries ----------

\newcommand{\RoleLineSize}{11.5pt}
\newcommand{\CompanyLineSize}{10.8pt}
\newcommand{\DurationLineSize}{9pt}
\newcommand{\BodyLineSize}{11pt}


% ---------------------------------------------------------------------------- %
% HEADER SPACING
% ---------------------------------------------------------------------------- %

% Space below name
\newcommand{\HeaderNameToKeyword}{12pt}

% Space between keyword row and contact row
\newcommand{\HeaderKeywordToContact}{5pt}

% Space between contact rows
\newcommand{\HeaderContactRows}{2pt}

% Space after entire header
\newcommand{\HeaderAfter}{-5pt}


% ---------------------------------------------------------------------------- %
% OPENING STATEMENT SPACING
% ---------------------------------------------------------------------------- %

% Space before opening statement
\newcommand{\StatementBefore}{0pt}

% Space after opening statement
\newcommand{\StatementAfter}{4pt}


% ---------------------------------------------------------------------------- %
% SECTION SPACING
% ---------------------------------------------------------------------------- %

% Space before a section
\newcommand{\SectionBefore}{4pt}

% Space between thick line and section title
\newcommand{\SectionLineToTitle}{0pt}

% Space between section title and thin line
\newcommand{\SectionTitleToLine}{2pt}

% Space between thin line and first entry
\newcommand{\SectionLineToEntry}{0pt}


% ---------------------------------------------------------------------------- %
% ENTRY SPACING
% ---------------------------------------------------------------------------- %

% Space before an entry
\newcommand{\EntryBefore}{0pt}

% Space between role and company
\newcommand{\RoleToCompany}{-5pt}

% Space between company and bullet list
\newcommand{\CompanyToBullets}{3pt}

% Space after an entry
\newcommand{\EntryAfter}{6pt}


% ---------------------------------------------------------------------------- %
% BULLET SPACING
% ---------------------------------------------------------------------------- %

% Space between individual bullets
\newcommand{\BulletSpacing}{2pt}

% Space before bullet list
\newcommand{\BulletBefore}{0pt}

% Space after bullet list
\newcommand{\BulletAfter}{0pt}

% Distance from left edge to bullet
\newcommand{\BulletIndent}{1.25em}


% ---------------------------------------------------------------------------- %
% HORIZONTAL LINES
% ---------------------------------------------------------------------------- %

% Thick section line
\newcommand{\SectionLineThickness}{1.5pt}

% Thin section line
\newcommand{\SubsectionLineThickness}{0.4pt}

% Main CV divider
\newcommand{\CVLineThickness}{0.5pt}


% ---------------------------------------------------------------------------- %
% DURATION COLUMN
% ---------------------------------------------------------------------------- %

% Width of the duration column
\newcommand{\DurationWidth}{4.5cm}


% ---------------------------------------------------------------------------- %
% ---------------------------------------------------------------------------- %
% END OF CONFIGURATION
% ---------------------------------------------------------------------------- %
% ---------------------------------------------------------------------------- %


% ---------------------------------------------------------------------------- %
% HYPERLINK CONFIGURATION
% ---------------------------------------------------------------------------- %

\hypersetup{
    colorlinks=true,
    urlcolor=CVBlack,
    linkcolor=CVBlack
}


% ---------------------------------------------------------------------------- %
% FONT SETUP
% ---------------------------------------------------------------------------- %

% The document-wide font family is selected above using \familydefault.


% ---------------------------------------------------------------------------- %
% GENERAL FORMATTING
% ---------------------------------------------------------------------------- %

\setlength{\parindent}{0pt}
\setlength{\parskip}{0pt}


% ---------------------------------------------------------------------------- %
% CV COMMANDS
% ---------------------------------------------------------------------------- %


% ------------------------------------------------------------
% Main horizontal divider
% ------------------------------------------------------------

\newcommand{\CVLine}{
    \vspace{\HeaderAfter}
    \noindent
    \textcolor{CVBlack}{
        \rule{\textwidth}{\CVLineThickness}
    }
    \vspace{\HeaderAfter}
}


% ------------------------------------------------------------
% First section
%
% The first section does NOT have a thick line because the
% CVLine separating the opening statement already provides it.
% ------------------------------------------------------------

\newcommand{\CVFirstSection}[1]{

    \vspace{\SectionBefore}

    % Section title

    {\fontsize{\SectionSize}{\SectionLineSize}
    \selectfont
    \textbf{\textcolor{CVBlack}{#1}}}

    \vspace{\SectionTitleToLine}

    % Thin grey line

    \noindent
    \textcolor{CVLightGray}{
        \rule{\textwidth}{\SubsectionLineThickness}
    }

    \vspace{\SectionLineToEntry}
}


% ------------------------------------------------------------
% Subsequent sections
%
% These have a thick black line above the section title.
% ------------------------------------------------------------

\newcommand{\CVSection}[1]{

    \vspace{\SectionBefore}

    % Thick black line

    \noindent
    \textcolor{CVBlack}{
        \rule{\textwidth}{\SectionLineThickness}
    }

    \vspace{\SectionLineToTitle}

    % Section title

    {\fontsize{\SectionSize}{\SectionLineSize}
    \selectfont
    \textbf{\textcolor{CVBlack}{#1}}}

    \vspace{\SectionTitleToLine}

    % Thin grey line

    \noindent
    \textcolor{CVLightGray}{
        \rule{\textwidth}{\SubsectionLineThickness}
    }

    \vspace{\SectionLineToEntry}
}


% ------------------------------------------------------------
% CV Entry
%
% #1 = Role / Degree
% #2 = Company / Institution
% #3 = Location
% #4 = Duration
% #5 = Bullet points
% ------------------------------------------------------------

\newcommand{\CVEntry}[5]{

    \vspace{\EntryBefore}

    % --------------------------------------------------------
    % Role + Duration
    % --------------------------------------------------------

    \noindent
    \begin{tabularx}{\textwidth}{
        @{}
        >{\raggedright\arraybackslash}X
        >{\raggedleft\arraybackslash}p{\DurationWidth}
        @{}
    }

        {\fontsize{\RoleSize}{\RoleLineSize}
        \selectfont
        \textbf{\textcolor{CVBlack}{#1}}}
        &
        {\fontsize{\DurationSize}{\DurationLineSize}
        \selectfont
        \textcolor{CVGray}{#4}}

    \end{tabularx}

    \vspace{\RoleToCompany}

    % --------------------------------------------------------
    % Company + Location
    % --------------------------------------------------------

    \noindent
    {\fontsize{\CompanySize}{\CompanyLineSize}
    \selectfont
    \textbf{\textcolor{CVDarkGray}{#2}}
    \textcolor{CVGray}{\textbar\ #3}}

    \vspace{\CompanyToBullets}

    % --------------------------------------------------------
    % Bullet points
    % --------------------------------------------------------

    \begin{itemize}[
        leftmargin=\BulletIndent,
        itemsep=\BulletSpacing,
        topsep=\BulletBefore,
        parsep=0pt,
        partopsep=0pt
    ]
        #5
    \end{itemize}

    \vspace{\EntryAfter}
}


% ------------------------------------------------------------
% Bullet point
% ------------------------------------------------------------

\newcommand{\CVBullet}[1]{
    \item
    {\fontsize{\BodySize}{\BodyLineSize}
    \selectfont
    \textcolor{CVDarkGray}{#1}}
}


% ---------------------------------------------------------------------------- %
% ---------------------------------------------------------------------------- %
% DOCUMENT
% ---------------------------------------------------------------------------- %
% ---------------------------------------------------------------------------- %

\begin{document}


% ---------------------------------------------------------------------------- %
% HEADER
% ---------------------------------------------------------------------------- %

\begin{center}

    % ------------------------------------------------------------------------ %
    % Name
    % ------------------------------------------------------------------------ %

    {
        \fontsize{\NameSize}{\NameLineSize}\selectfont
        \textbf{\textcolor{CVBlack}{Jason Wakim Richards}}
    }

    % ------------------------------------------------------------------------ %
    % Contact information
    % ------------------------------------------------------------------------ %

    \vspace{\HeaderKeywordToContact}

    {
        \fontsize{\ContactSize}{\ContactLineSize}
        \selectfont
        \textcolor{CVBlack}
        {
            +44 77 08 348 077
            \quad\textbar\quad
            \href{mailto:jason2richards@gmail.com}{jason2richards@gmail.com}
            \\
            \href{https://github.com/richarJ2002}{github.com/richarJ2002}
            \quad\textbar\quad
            \href{https://www.linkedin.com/in/jason-richards-jwr2002/}
            {linkedin.com/in/jason-richards-jwr2002/}
            \quad\textbar\quad
            \href{https://richarj2002.github.io/personal_website/}{richarj2002.github.io/personal\_website/}
        }
    }

\end{center}


% ---------------------------------------------------------------------------- %
% OPENING STATEMENT
% ---------------------------------------------------------------------------- %

\CVLine

\vspace{\StatementBefore}

{
    \fontsize
    {\StatementSize}
    {\StatementLineSize}
    \selectfont
    \textcolor{CVDarkGray}
    {
        Guidance, navigation, and control engineer with hands-on experience developing,
        integrating, and testing autonomy for planetary rovers, spacecraft, and UAVs.
        At Airbus Defence \& Space, I led LiDAR test activities for an autonomous Sample
        Fetch Rover concept, built a networked sensor test bench, characterized sensor
        performance, and developed C/C++ simulation software that generated synthetic
        imagery and LiDAR scans. I now develop and flight-test UAV localization,
        mapping, planning, and control using ROS 2, PX4, Linux, Python, C++, cameras,
        and IMUs. As ADCS Lead for the JamSail 3U CubeSat, I defined requirements and
        interfaces, developed embedded flight software and simulation, and implemented
        an EKF for multi-sensor attitude estimation. I bring practical experience
        across GNC requirements, algorithm development, sensor calibration, performance
        analysis, hardware/software integration, and verification and validation.
    }
}

\vspace{\StatementAfter}

\CVLine


% ---------------------------------------------------------------------------- %
% WORK EXPERIENCE
% --
\CVFirstSection{Work Experience}

\CVEntry
{Doctoral Researcher -- Autonomous UAV Systems}
{SnT, University of Luxembourg}
{Kirchberg, Luxembourg}
{November 2025 -- Present}
{
    \CVBullet{Develop C++ and Python localization, Semantic Relational SLAM, and situational-awareness algorithms for UAVs, linking camera/IMU perception and state estimation with planning and control.}

    \CVBullet{Integrate, tune, and deploy autonomous navigation on physical UAVs using ROS 2, PX4, embedded Linux, Raspberry Pi, and OpenVINS visual--inertial odometry.}

    \CVBullet{Build reproducible Docker-based development and deployment environments and camera--IMU characterization and calibration workflows with Kalibr, supported by version-controlled software and technical documentation.}

    \CVBullet{Plan and conduct simulation, laboratory, hardware, and flight testing; analyze integrated performance and troubleshoot interfaces across sensing, estimation, mapping, planning, and control.}

    \CVBullet{Collaborate across perception, autonomy, and vehicle-integration activities while documenting algorithms, interfaces, test results, and technically grounded next steps.}
}

\CVEntry
{AOCS/GNC Engineer Intern}
{Airbus Defence \& Space}
{Stevenage, UK}
{July 2023 -- August 2024}
{
    \CVBullet{Served as lead test engineer for LiDAR activities during Integrated BreadBoarding (iBB) 3 \& 4, planning procedures and supporting end-to-end verification of the Sample Fetch Rover concept at Technology Readiness Level 6.}

    \CVBullet{Designed and built an integrated hardware/software GNC test bench, connecting LiDAR to rover breadboards through ROS over a local wireless network to capture data, characterize sensor performance, and validate navigation models.}

    \CVBullet{Developed C/C++ software on Linux for an offline GNC simulation, coupling a natural-scene generator to produce synthetic camera imagery and LiDAR scans for rover algorithm development and performance analysis.}

    \CVBullet{Investigated LiDAR and radar simulation, terrain generation, and Vector Field Histogram planning for lunar-rover navigation, translating prototype concepts into the existing simulation environment.}

    \CVBullet{Analyzed test-campaign data and documented results while working across GNC, simulation, software, and test disciplines using Git, GitLab, Jira, and Confluence in a fast-paced Agile team.}
}

% ---------------------------------------------------------------------------- %
% PROJECTS & LEADERSHIP
% ---------------------------------------------------------------------------- %

\CVSection{Projects \& Leadership}

\CVEntry
{JamSail 3U CubeSat -- AOCS/ADCS Lead}
{NottSpace, University of Nottingham}
{Nottingham, UK}
{September 2024 -- August 2025}
{
    \CVBullet{Led ADCS design and implementation for a NottSpace 3U CubeSat, defining subsystem and algorithm requirements, concept of operations, architecture, interfaces, and a verification and validation plan.}

    \CVBullet{Developed embedded C flight software and a C/C++/CMake GNC simulation on Linux to design and verify attitude determination and control algorithms before deployment to flight hardware.}

    \CVBullet{Designed and implemented an EKF that fuses sun-sensor, magnetometer, and IMU measurements, together with TRIAD attitude determination and control-state logic.}

    \CVBullet{Characterized sensor noise with FFT analysis, tuned digital low-pass filters, and verified their performance through simulation and hardware testing; implemented the IGRF magnetic-field model as a physics-based simulation input.}
}

% Start the remaining projects on page two, keeping the complete JamSail entry
% on page one and preserving Education & Training as a single final section.
\clearpage

\CVEntry
{IDET-Sat -- Systems Engineer}
{Nottingham Space Society}
{Nottingham, UK}
{September 2022 -- July 2023}
{
    \CVBullet{Defined system requirements, architecture, subsystem interfaces, Project Breakdown Tree, and Work Breakdown Structure within a student-led, 12-person team developing a 3U CubeSat.}

    \CVBullet{Developed software on a Raspberry Pi running Linux and integrated a microcontroller, IMU, and cameras over I2C before the concept progressed to the competition build phase.}

    \CVBullet{Contributed across embedded software, hardware and sensor integration, and system simulation while coordinating interfaces and technical development within the student team.}

    \CVBullet{Achieved a top-four finish among 12 UK teams as first-time entrants to the UKSEDS Satellite Design Competition.}

    \CVBullet{Led cross-team meetings, communicated design changes, and presented the integrated concept to industry experts, balancing technical interfaces, build constraints, and competition deadlines.}
}

\CVEntry
{UKSEDS Volunteer}
{UK Students for the Exploration and Development of Space}
{Remote \& UK field events}
{Sept. 2024 -- Sept. 2026}
{
    \CVBullet{Served as an Olympus Rover Trials (ORT) volunteer in the first year, helping organise a UK national rover competition with Airbus Defence \& Space and RAL Space.}

    \CVBullet{Progressed in the second year to Competitions Management Co-Lead, overseeing four national competitions: ORT, SDC, ISAM, and NRC.}

    \CVBullet{Coordinated volunteers, schedules, technical milestones, participant support, and partner engagement with Airbus Defence \& Space (SDC), the UK Space Agency (ISAM), Midlands Rocketry Club, and Skyrora (NRC).}

    \CVBullet{Produced clear participant documentation and supported venue and field-event logistics, helping multidisciplinary volunteer teams deliver competitions to fixed schedules.}
}

% ---------------------------------------------------------------------------- %
% EDUCATION
% ---------------------------------------------------------------------------- %

\CVSection{Education \& Training}

\CVEntry
{MEng Aerospace Engineering -- First-Class Honours}
{University of Nottingham}
{Nottingham, UK}
{September 2020 -- August 2025}
{
    \CVBullet{Graduated with First-Class Honours, specializing in Aerospace Control Systems, Avionics, dynamics, software, modeling, and simulation.}

    \CVBullet{Applied C/C++, Python, MATLAB, Simulink, and Linux to estimation, control, flight dynamics, and autonomous-system projects.}

    \CVBullet{Completed design projects spanning aerospace systems, avionics, simulation, and control, providing the systems-engineering foundation for later space and robotics work.}
}

\CVEntry
{IEEE RAS Summer School on Multi-Robot Systems}
{Czech Technical University in Prague}
{Prague, Czechia}
{July 2026}
{
    \CVBullet{Completed the 2026 IEEE RAS Summer School on Multi-Robot Systems, covering cooperative aerial robotics and active perception; presented the SaCoGraph collaborative-UAV research concept.}

    \CVBullet{Applied multi-robot autonomy concepts to collaborative sensing, semantic mapping, and resource-aware mission planning.}
}

\CVEntry
{Space Resources \& Robotics Summer School}
{AGH University of Kraków}
{Krakow, Poland}
{May 2026}
{
    \CVBullet{Completed intensive training in space resources, planetary robotics, and robotic manipulation for exploration applications.}

    \CVBullet{Developed a broader systems perspective on robotic operations in challenging off-world environments.}
}

\CVEntry
{Space Systems Engineering Training Course}
{ESA Academy}
{Transinne, Belgium}
{May 2025}
{
    \CVBullet{Completed intensive ESA Academy training across the space-system lifecycle, from mission objectives and concepts of operations through architecture, integration, verification, and operations.}

    \CVBullet{Applied requirements engineering and concurrent-design methods to translate user needs into traceable, testable requirements and manage interfaces between technical disciplines.}

    \CVBullet{Practised system options and trade studies, budgets and margins, verification and validation planning, and the relationship between engineering decisions, schedule, cost, and risk.}

    \CVBullet{Completed multidisciplinary group exercises and communicated design rationale and systems-engineering deliverables, strengthening collaborative problem-solving across complex spacecraft projects.}
}


\end{document}